\chapter{Sintonización PID mediante Grey Wolf Optimizer (GWO)}
\label{ch:gwo}

\section{Introducción del capítulo}
El Grey Wolf Optimizer pertenece a la familia de algoritmos de enjambre y ha ganado presencia
en problemas de control por su estructura sencilla y su buena capacidad para balancear
exploración y explotación. Aunque el borrador principal del proyecto prioriza PSO, la inclusión
de GWO en este libro permite ampliar la comparación y examinar si una estrategia de jerarquía
social distinta puede producir sintonías competitivas para el mismo problema hidráulico.

\section{Objetivo del capítulo}
Este capítulo presenta el fundamento del GWO y su adaptación al ajuste de un controlador PID.
La meta es mostrar cómo la metáfora de caza cooperativa puede traducirse en una regla de
actualización útil para recorrer el espacio de ganancias del controlador bajo restricciones
análogas a las empleadas con PSO.

\section{Inspiración del algoritmo}
GWO se inspira en la jerarquía de una manada de lobos grises. Los individuos alfa, beta y
delta representan las mejores soluciones conocidas y guían el movimiento del resto del grupo,
denominado omega. En términos de optimización, esta jerarquía introduce una forma de
liderazgo distribuido: la población no sigue a un único óptimo provisional, sino a una tríada de
referencias que ayuda a reducir el riesgo de convergencia prematura.

\section{Modelo matemático del GWO}
La posición de cada solución se actualiza a partir de su distancia relativa respecto de los tres
mejores lobos del grupo. De forma resumida, el algoritmo calcula posiciones intermedias
asociadas a alfa, beta y delta, y luego combina esas referencias para obtener la nueva solución
candidata. El parámetro $a$, decreciente con las iteraciones, controla la transición desde una
fase más exploratoria hacia una fase de refinamiento, lo que constituye uno de los mecanismos
clásicos de convergencia del método.

\section{Pseudocódigo}
\begin{algorithm}[H]
\caption{Esquema general del GWO para sintonización PID}
\begin{algorithmic}[1]
\State Inicializar una población de soluciones candidatas
\State Evaluar cada lobo con la función de costo del problema
\State Identificar las soluciones alfa, beta y delta
\While{no se cumpla el criterio de parada}
    \For{cada lobo}
        \State Calcular su distancia respecto de alfa, beta y delta
        \State Actualizar la posición usando la combinación de esas tres referencias
        \State Restringir la solución al dominio admisible
    \EndFor
    \State Revaluar la población y actualizar alfa, beta y delta
\EndWhile
\State Entregar la mejor sintonización encontrada
\end{algorithmic}
\end{algorithm}

\section{Parámetros de ajuste}
Los parámetros más influyentes del GWO son el tamaño de la población, el número de
iteraciones y la evolución del coeficiente de contracción $a$. A diferencia de otros métodos, el
algoritmo no requiere demasiados hiperparámetros, lo que simplifica su uso comparativo. Esta
economía de configuración es una de las razones por las que se ha convertido en un candidato
frecuente para tareas de sintonización en sistemas eléctricos y de potencia.

\section{Aplicación a la sintonización PID}
En la sintonización PID, cada lobo codifica una tripleta $\left[K_p,K_i,K_d\right]$ y su calidad se
evalúa mediante los mismos escenarios de simulación empleados con PSO. Esta uniformidad
metodológica es importante porque permite comparar algoritmos sin introducir sesgos en la
planta, la función objetivo o las perturbaciones. El interés principal consiste en observar si la
dinámica de liderazgo múltiple del GWO ofrece ventajas en robustez o en calidad final de la
respuesta de frecuencia.

\section{Cierre del capítulo}
GWO aporta una alternativa sobria y competitiva dentro de la familia de métodos de enjambre.
Su valor dentro de este libro radica en ampliar el horizonte comparativo más allá del PSO. El
siguiente capítulo introduce Eagle Strategy, una propuesta que combina búsqueda global y
refinamiento local con una lógica diferente de movimiento.
